% 2007 年重庆卷文科第 21 题：抛物线 y^2=8x、焦点 F、准线 l 与过 F 的直线交于 A、B；
%   弦 AB 的垂直平分线 m 交 x 轴于 P，倾斜角 α 标在直线与 x 轴之间。
% 来源 PDF：gaokaomath-src/普通高考/2007/2007重庆文.pdf
% SHA-256：a647b60227db10d70c5c3fbe72b6ffe2c8a0667b233695e55bb4f5a7cef633a6
% 原页：2007重庆文.pdf 第 2 页（A3 三栏）第 2 栏，书页码 358。
% 重绘依据：按 α=60° 作图（A 在上、B 在下、P 在焦点右侧）；m 取过弦中点 M(10/3,2√3)、
%   斜率 −1/√3 的直线，与 x 轴的交点即 P(22/3,0)，与原页位置关系一致。
\begin{tikzpicture}[x=0.78cm, y=0.52cm, semithick,
  every node/.style={inner sep=1.5pt}]

  % 坐标轴
  \draw[->,>=stealth] (-2.8,0) -- (8.6,0) node[right] {$x$};
  \draw[->,>=stealth] (0,-7.6) -- (0,8.2) node[above] {$y$};
  \node[below left] at (0,0) {$O$};

  % 准线 l：x=-2
  \draw (-2,-7.2) -- (-2,7.6) node[above] {$l$};

  % 抛物线 y^2=8x
  \draw[domain=-7.2:7.2,samples=160,line width=1pt] plot ({\x*\x/8},{\x});

  % 过 F 的直线（α=60°），交抛物线于 A、B
  \coordinate (F) at (2,0);
  \coordinate (A) at (6,6.928);
  \coordinate (B) at (0.667,-2.309);
  \draw (B) -- (A);

  % 弦 AB 的垂直平分线 m（过中点 M，斜率 -1/√3），交 x 轴于 P
  \coordinate (P) at (7.333,0);
  \draw (1.2,3.539) -- (8.3,-0.559);
  \node[above left] at (1.2,3.539) {$m$};

  % 角 α
  \draw (2.9,0) arc[start angle=0,end angle=60,radius=0.9];
  \node at (3.35,0.95) {$\alpha$};

  \fill (F) circle (1.3pt);
  \fill (P) circle (1.3pt);
  \node[below right=2pt] at (F) {$F$};
  \node[below=3pt] at (P) {$P$};
  \node[above left] at (A) {$A$};
  \node[below left] at (B) {$B$};
\end{tikzpicture}
